\section{Codifica vocale}
\subsection{Caratteristiche della voce e classi di codificatori}
\subsubsection*{Suoni vocali e non vocali}
Il segnale vocale può essere considerato un segnale localmente stazionario su periodi di 20 ms. In tali periodi, il carattere
del segnale rimane costante per la stazionarietà durante la pronuncia di un fonema. In generale i periodi di stazionarietà
si classificano in due categorie principali:
\begin{itemize}
	\item suoni vocali: segnali pseudoperiodici, associati alle vocali e alle consonanti s e f, caratterizzati da alcune
	frequenze predominanti dovute alla periodicità data dalla vibrazione delle corde vocali
	\item suoni non vocali: segnali aperiodici, associati alle consonanti p, t, k, s, f, h, caratterizzati da un
	spettro piuttosto uniforme, simile al rumore bianco
	\item suoni misti: transizioni tra i diversi fonemi
\end{itemize}

\subsubsection*{Classi di codificatori}
\begin{itemize}
	\item \textbf{codificatori a forme d'onda}: codificano la forma d'onda del segnale, senza ottimizzazioni per le
	caratteristiche del segnale vocale; hanno bassa complessità, alta qualità del parlato, ma bitrate elevato; un esempio
	è la codifica PCM
	\item \textbf{codificatori vocali o vocoder}: modellano il modo in cui il segnale vocale viene generato, stimando
	i parametri del modello e trasmettendo solo questi parametri; hanno bassa qualità del parlato, ma bitrate molto basso;
	un esempio è la codifica LPC-10
	\item \textbf{codificatori ibridi}: combinano le due tecniche precedenti, utilizzano tecniche di vocoding unite a
	codifiche a forme d'onda per il residui di predizione; hanno qualità del parlato migliore dei vocoder e bitrate più
	elevato; esempi sono ACELP, AMR-WB e Opus
\end{itemize}

\subsection{Codificatore vocoder basato su LPC-10}
\subsubsection*{Codifica a predizione lineare (LPC-10)}
La codifica a predizione lineare LPC-10 consiste nel produrre una predizione \(x_p(n)\) del campione corrente \(x(n)\)
attraverso una combinazione lineare dei 10 campioni precedenti \(x(n\!-\!1), \; x(n\!-\!2), \; \dots \; x(n\!-\!10)\)
pesati da opportuni coefficienti \(a_1, a_2, \dots, a_{10}\).
\[x_p(n) = \sum_{k=1}^{10} a_k \, x(n-k)\]
I coefficienti \(a_k\) vengono calcolati minimizzando l'errore di predizione quadratico medio tra il campione
corrente \(x(n)\) e la sua predizione \(x_p(n)\). L'errore di predizione \(y(n)\) è definito come segue:
\[y(n) = x(n) - x_p(n) = x(n) - a_1 \, x(n-1) - a_2 \, x(n-2) - \dots - a_{10} \, x(n-10)\]
Si osserva che conoscendo i coefficienti \(a_k\) e l'errore di predizione \(y(n)\), è possibile ricostruire esattamente
il campione corrente \(x(n)\). Risulta sufficiente che il codificatore trasmetta tali informazioni al decodificatore
(opportunamente quantizzate) per permettere la ricostruzione del segnale vocale.

\subsubsection*{Trasmissione del residuo}
Se i coefficienti \(a_k\) sono stati calcolati correttamente, il residuo \(y(n)\) risulta essere:
\begin{itemize}
	\item un fruscio o rumore bianco per i suoni non vocali, di cui viene trasmesso solo la potenza, siccome il
	decodificatore è in grado di generare da sé un rumore bianco casuale
	\item un segnale pseudoperiodico con frequenza predominante detta pitch per i suoni vocali, che viene trasmesso
	al decodificatore come treno di impulsi periodici
\end{itemize}
Per individuare pitch e potenza del residuo si utilizzano tecniche come la stima dell'autocorrelazione.

\subsubsection*{Tasso di codifica}
Supponendo di suddividere il segnale vocale in finestre di 20 ms, per ogni finestra si trasmettono:
\begin{itemize}
	\item 36 bit per i 10 coefficienti del filtro di predizione lineare
	\item 1 bit per distinguere se l'errore di predizione è un treno di impulsi o rumore bianco
	\item 6 bit per la potenza del residuo, sia per il rumore bianco che per il treno di impulsi
	\item 7 bit per il pitch del residuo, solo se l'errore di predizione è un treno di impulsi
\end{itemize}
Complessivamente si ottiene un tasso di codifica medio di 2.4 kbps, che è molto più basso rispetto ai 64 kbps della
codifica PCM, a discapito di una qualità del parlato molto più bassa.
\[R_\text{voiced} = \frac{36 + 1 + 6 + 7 \; \text{bit}}{20 \; \text{ms}} = 2500 \; \text{bps} \qquad R_\text{unvoiced} = \frac{36 + 1 + 6 \; \text{bit}}{20 \; \text{ms}} = 2150 \; \text{bps} \qquad R_\text{avg} = 2400 \; \text{bps}\]

\subsubsection*{Criticità del codificatore LPC-10}
Il codificatore LPC-10 ha alcune criticità che lo rendono poco utilizzabile in applicazioni pratiche:
\begin{itemize}
	\item \textbf{decisione binaria}: la classificazione dei campioni in ``vocali'' e ``non vocali'' risulta troppo
	stringente in quanto non tiene conto degli stati intermedi di transizione tra i fonemi; sarebbe opportuno
	avere a disposizione più classi di campioni per catturare meglio la variabilità del segnale vocale
	\item \textbf{errori di classificazione}: se il codificatore commette errori nel classificare un campione, si
	verificheranno distorsioni e perdita di informazioni nel segnale ricostruito creando un effetto ``sussurrato''
	o ``soffiato'' nel parlato; per risolvere ciò si possono utilizzare modelli con catene di feedback
	\item \textbf{semplificazione del residuo}: trattando il residuo come generico rumore bianco, si perdono informazioni
	importanti sulla struttura spettrale e sulla coerenza di fase del segnale originale, producendo un timbro artificiale;
	risulta utile avere più parametri per modellare meglio il residuo
	\item \textbf{latenza}: i codificatori a predizione lineare possiedono una certa latenza dovuta in parte al costo
	computazionale, in parte ad eventuali buffer necessari nelle implementazioni pratiche
	\item \textbf{resilienza}: i codificatori a predizione lineare sono molto sensibili agli errori casuali e agli errori
	a burst, per cui è necessario impiegare opportune tecniche di FEC (Forward Error Correction)
	\item \textbf{soppressione del silenzio}: per risparmiare banda, si può aggiungere un modulo di rilevamento dell'attività
	vocale (VAD, Voice Activity Detection) che interrompe la trasmissione del segnale durante i periodi di silenzio, ma allo
	stesso tempo può portare a improvvise interruzioni del parlato
\end{itemize}

\subsection{Codificatori più avanzati - ACELP, AMR-WB e Opus}
\subsubsection*{Codificatore ACELP - G.729}
Il codificatore ACELP (Algebraic Code-Excited Linear Prediction) è un codificatore basato su predizione lineare che
dispone di un set molto esteso di residui (generati da codici algebrici) per eccitare il filtro di predizione. Ciò
permette di modellare meglio gli stati intermedi di transizione tra fonemi e dispone di feedback loop per migliorare
le prestazioni. Possiede anche moduli per la gestione dei mascheramenti.

\subsubsection*{Codificatore AMR-WB - G.722.2}
Il codificatore Adaptive Multi-Rate Wideband (AMR-WB) è un codificatore ibrido utilizzato come standard per l'HD-Voice
nelle reti mobili LTE/5G. La particolarità di questo codificatore è l'estensione della banda passante fino a 7 kHz (contro
i 4 kHz della banda telefonica tradizionale) e la possibilità di variare il bitrate da 6.6 kbps a 23.85 kbps in base alla
qualità del canale radio. Questo permette maggiore qualità del parlato grazie all'aumento del contenuto armonico delle alte
frequenze e una maggiore robustezza della trasmissione in caso di canali radio con banda limitata.

\newpage

\subsection{Codificatore Opus - RFC 6716}
\subsubsection*{Caratteristiche generali}
Il codificatore Opus è un codificatore ibrido con bitrate variabile da 6 kbps a 510 kbps e banda estesa da 8 kHz a 48 kHz,
ampiamente utilizzato per servizi di videochiamata WebRTC (meet, zoom, teams), messaggistica (whatsapp, telegram), streaming
audio (youtube, spotify) e gaming (discord, twitch).

\subsubsection*{Architettura ibrida SILT + CELT}
Opus si basa su due motori di elaborazione paralleli per codificare il segnale vocale e il segnale musicale:
\begin{itemize}
	\item \textbf{motore SILT}: utilizza la predizione lineare (LPC) basata sul modello del tratto vocale umano, come visto
	in precedenza; si occupa della parte inferiore dello spettro armonico (\(\leq 8 \; \text{kHz}\)) tipica del parlato
	\item \textbf{motore CELT}: utilizza una variante ottimizzata della trasformata discreta di Fourier (MDCT), simile
	a quella utilizzata nei codificatori MP3 e AAC; si occupa di codificare la parte superiore dello spettro armonico
	(\(> 8 \; \text{kHz}\)) tipica della musica e in generale dei suoni non vocali
\end{itemize}

\noindent
Se il segnale è prevalentemente vocale, il codificatore utilizza solo il motore SILT, mentre se il segnale contiene anche
componenti musicali, il codificatore utilizza entrambi i motori in parallelo. In questo modo, si riesce a garantire una
qualità del parlato molto elevata con bassi bitrate e contemporaneamente una qualità musicale paragonabile a quella dei
codificatori MP3 e AAC.

\subsubsection*{Diffusione e vantaggi su altri codificatori}
La sua grande diffusione è dovuta al fatto che ha una bassa latenza algoritmica e strutturale di soli 26.5 ms, dispone di
grande flessibilità nella gestione del bitrate e ha una banda sufficientemente estesa sia per segnali vocali che musicali.
Inoltre è open source e royalty-free, per cui può essere utilizzato liberamente senza costi di licenza da chiunque necessiti
di un codificatore universale di qualità.
